% =====================================================================
% Quiz 1 -- Cálculo 11° -- Trimestre I -- Semana 02 (sesión 004, jue 13 ago.)
% Clave:  pdflatex -jobname=quiz-clave "\def\clave{}% =====================================================================
% Quiz 1 -- Cálculo 11° -- Trimestre I -- Semana 02 (sesión 004, jue 13 ago.)
% Clave:  pdflatex -jobname=quiz-clave "\def\clave{}% =====================================================================
% Quiz 1 -- Cálculo 11° -- Trimestre I -- Semana 02 (sesión 004, jue 13 ago.)
% Clave:  pdflatex -jobname=quiz-clave "\def\clave{}% =====================================================================
% Quiz 1 -- Cálculo 11° -- Trimestre I -- Semana 02 (sesión 004, jue 13 ago.)
% Clave:  pdflatex -jobname=quiz-clave "\def\clave{}\input{quiz}"  (dos veces)
% =====================================================================
\documentclass[11pt]{article}
\makeatletter\def\input@path{{../../../../../plantillas/estilo/}}\makeatother
\usepackage{mmcantillo}

\tipodocumento{Quiz}
\asignatura{Cálculo}
\titulo{Quiz 1: los sistemas numéricos}
\grado{11°}
\periodo{I}
\tiempo{20 min}
% DBA: matematicas grado 11 · DBA 1 — Utiliza las propiedades de los números (naturales,
%      enteros, racionales y reales) y sus relaciones y operaciones para construir y comparar
%      los distintos sistemas numéricos.
% Evidencias: (1) propiedades comunes y diferentes (puntos 1 y 4); comparar sistemas
%      numéricos (punto 2); representaciones (punto 3).

\begin{document}

\encabezado

\begin{instrucciones}
Trabajo individual. Escribe los procedimientos: las respuestas sin justificación no tienen
puntos.
\end{instrucciones}

\begin{preguntas}

\pregunta[1] ¿Cuál de las siguientes afirmaciones es \textbf{falsa}?
\begin{opciones*}(2)
  \item Todo número natural es entero.
  \item Todo número entero es racional.
  \item Todo número racional es entero.
  \item El $0$ es un número natural.
\end{opciones*}
\begin{solucion}
c). Contraejemplo: $\frac{1}{2} \in \Q$ pero $\frac{1}{2} \notin \Z$.
\end{solucion}

\pregunta[1.5] Ordena de menor a mayor y explica cómo los comparaste:
\[
  -\frac{3}{4}, \qquad 0{,}7, \qquad -0{,}8, \qquad \frac{2}{3}
\]
\begin{solucion}[3cm]
En decimales: $-0{,}75$; $0{,}7$; $-0{,}8$; $0{,}\overline{6}$. Orden:
$-0{,}8 < -\frac{3}{4} < \frac{2}{3} < 0{,}7$.
\end{solucion}

\pregunta[1] Escribe $0{,}\overline{36}$ como fracción irreducible.
\begin{solucion}[2.5cm]
$100x - x = 36$, $x = \frac{36}{99} = \frac{4}{11}$.
\end{solucion}

\pregunta[1.5] Di si cada afirmación es verdadera o falsa. Si es falsa, da un contraejemplo;
si es verdadera, explica por qué.
\begin{opciones}
  \item La resta de números enteros es conmutativa.
  \item Todo número racional distinto de cero tiene inverso multiplicativo en $\Q$.
  \item El conjunto $\N$ es cerrado para la resta.
\end{opciones}
\begin{solucion}[2.5cm]
a) Falsa: $5 - 3 = 2$ pero $3 - 5 = -2$. b) Verdadera: el inverso de $\frac{p}{q}$ (con
$p \neq 0$) es $\frac{q}{p}$, que también es racional. c) Falsa: $3 - 5 = -2 \notin \N$.
\end{solucion}

\end{preguntas}

\end{document}
"  (dos veces)
% =====================================================================
\documentclass[11pt]{article}
\makeatletter\def\input@path{{../../../../../plantillas/estilo/}}\makeatother
\usepackage{mmcantillo}

\tipodocumento{Quiz}
\asignatura{Cálculo}
\titulo{Quiz 1: los sistemas numéricos}
\grado{11°}
\periodo{I}
\tiempo{20 min}
% DBA: matematicas grado 11 · DBA 1 — Utiliza las propiedades de los números (naturales,
%      enteros, racionales y reales) y sus relaciones y operaciones para construir y comparar
%      los distintos sistemas numéricos.
% Evidencias: (1) propiedades comunes y diferentes (puntos 1 y 4); comparar sistemas
%      numéricos (punto 2); representaciones (punto 3).

\begin{document}

\encabezado

\begin{instrucciones}
Trabajo individual. Escribe los procedimientos: las respuestas sin justificación no tienen
puntos.
\end{instrucciones}

\begin{preguntas}

\pregunta[1] ¿Cuál de las siguientes afirmaciones es \textbf{falsa}?
\begin{opciones*}(2)
  \item Todo número natural es entero.
  \item Todo número entero es racional.
  \item Todo número racional es entero.
  \item El $0$ es un número natural.
\end{opciones*}
\begin{solucion}
c). Contraejemplo: $\frac{1}{2} \in \Q$ pero $\frac{1}{2} \notin \Z$.
\end{solucion}

\pregunta[1.5] Ordena de menor a mayor y explica cómo los comparaste:
\[
  -\frac{3}{4}, \qquad 0{,}7, \qquad -0{,}8, \qquad \frac{2}{3}
\]
\begin{solucion}[3cm]
En decimales: $-0{,}75$; $0{,}7$; $-0{,}8$; $0{,}\overline{6}$. Orden:
$-0{,}8 < -\frac{3}{4} < \frac{2}{3} < 0{,}7$.
\end{solucion}

\pregunta[1] Escribe $0{,}\overline{36}$ como fracción irreducible.
\begin{solucion}[2.5cm]
$100x - x = 36$, $x = \frac{36}{99} = \frac{4}{11}$.
\end{solucion}

\pregunta[1.5] Di si cada afirmación es verdadera o falsa. Si es falsa, da un contraejemplo;
si es verdadera, explica por qué.
\begin{opciones}
  \item La resta de números enteros es conmutativa.
  \item Todo número racional distinto de cero tiene inverso multiplicativo en $\Q$.
  \item El conjunto $\N$ es cerrado para la resta.
\end{opciones}
\begin{solucion}[2.5cm]
a) Falsa: $5 - 3 = 2$ pero $3 - 5 = -2$. b) Verdadera: el inverso de $\frac{p}{q}$ (con
$p \neq 0$) es $\frac{q}{p}$, que también es racional. c) Falsa: $3 - 5 = -2 \notin \N$.
\end{solucion}

\end{preguntas}

\end{document}
"  (dos veces)
% =====================================================================
\documentclass[11pt]{article}
\makeatletter\def\input@path{{../../../../../plantillas/estilo/}}\makeatother
\usepackage{mmcantillo}

\tipodocumento{Quiz}
\asignatura{Cálculo}
\titulo{Quiz 1: los sistemas numéricos}
\grado{11°}
\periodo{I}
\tiempo{20 min}
% DBA: matematicas grado 11 · DBA 1 — Utiliza las propiedades de los números (naturales,
%      enteros, racionales y reales) y sus relaciones y operaciones para construir y comparar
%      los distintos sistemas numéricos.
% Evidencias: (1) propiedades comunes y diferentes (puntos 1 y 4); comparar sistemas
%      numéricos (punto 2); representaciones (punto 3).

\begin{document}

\encabezado

\begin{instrucciones}
Trabajo individual. Escribe los procedimientos: las respuestas sin justificación no tienen
puntos.
\end{instrucciones}

\begin{preguntas}

\pregunta[1] ¿Cuál de las siguientes afirmaciones es \textbf{falsa}?
\begin{opciones*}(2)
  \item Todo número natural es entero.
  \item Todo número entero es racional.
  \item Todo número racional es entero.
  \item El $0$ es un número natural.
\end{opciones*}
\begin{solucion}
c). Contraejemplo: $\frac{1}{2} \in \Q$ pero $\frac{1}{2} \notin \Z$.
\end{solucion}

\pregunta[1.5] Ordena de menor a mayor y explica cómo los comparaste:
\[
  -\frac{3}{4}, \qquad 0{,}7, \qquad -0{,}8, \qquad \frac{2}{3}
\]
\begin{solucion}[3cm]
En decimales: $-0{,}75$; $0{,}7$; $-0{,}8$; $0{,}\overline{6}$. Orden:
$-0{,}8 < -\frac{3}{4} < \frac{2}{3} < 0{,}7$.
\end{solucion}

\pregunta[1] Escribe $0{,}\overline{36}$ como fracción irreducible.
\begin{solucion}[2.5cm]
$100x - x = 36$, $x = \frac{36}{99} = \frac{4}{11}$.
\end{solucion}

\pregunta[1.5] Di si cada afirmación es verdadera o falsa. Si es falsa, da un contraejemplo;
si es verdadera, explica por qué.
\begin{opciones}
  \item La resta de números enteros es conmutativa.
  \item Todo número racional distinto de cero tiene inverso multiplicativo en $\Q$.
  \item El conjunto $\N$ es cerrado para la resta.
\end{opciones}
\begin{solucion}[2.5cm]
a) Falsa: $5 - 3 = 2$ pero $3 - 5 = -2$. b) Verdadera: el inverso de $\frac{p}{q}$ (con
$p \neq 0$) es $\frac{q}{p}$, que también es racional. c) Falsa: $3 - 5 = -2 \notin \N$.
\end{solucion}

\end{preguntas}

\end{document}
"  (dos veces)
% =====================================================================
\documentclass[11pt]{article}
\makeatletter\def\input@path{{../../../../../plantillas/estilo/}}\makeatother
\usepackage{mmcantillo}

\tipodocumento{Quiz}
\asignatura{Cálculo}
\titulo{Quiz 1: los sistemas numéricos}
\grado{11°}
\periodo{I}
\tiempo{20 min}
% DBA: matematicas grado 11 · DBA 1 — Utiliza las propiedades de los números (naturales,
%      enteros, racionales y reales) y sus relaciones y operaciones para construir y comparar
%      los distintos sistemas numéricos.
% Evidencias: (1) propiedades comunes y diferentes (puntos 1 y 4); comparar sistemas
%      numéricos (punto 2); representaciones (punto 3).

\begin{document}

\encabezado

\begin{instrucciones}
Trabajo individual. Escribe los procedimientos: las respuestas sin justificación no tienen
puntos.
\end{instrucciones}

\begin{preguntas}

\pregunta[1] ¿Cuál de las siguientes afirmaciones es \textbf{falsa}?
\begin{opciones*}(2)
  \item Todo número natural es entero.
  \item Todo número entero es racional.
  \item Todo número racional es entero.
  \item El $0$ es un número natural.
\end{opciones*}
\begin{solucion}
c). Contraejemplo: $\frac{1}{2} \in \Q$ pero $\frac{1}{2} \notin \Z$.
\end{solucion}

\pregunta[1.5] Ordena de menor a mayor y explica cómo los comparaste:
\[
  -\frac{3}{4}, \qquad 0{,}7, \qquad -0{,}8, \qquad \frac{2}{3}
\]
\begin{solucion}[3cm]
En decimales: $-0{,}75$; $0{,}7$; $-0{,}8$; $0{,}\overline{6}$. Orden:
$-0{,}8 < -\frac{3}{4} < \frac{2}{3} < 0{,}7$.
\end{solucion}

\pregunta[1] Escribe $0{,}\overline{36}$ como fracción irreducible.
\begin{solucion}[2.5cm]
$100x - x = 36$, $x = \frac{36}{99} = \frac{4}{11}$.
\end{solucion}

\pregunta[1.5] Di si cada afirmación es verdadera o falsa. Si es falsa, da un contraejemplo;
si es verdadera, explica por qué.
\begin{opciones}
  \item La resta de números enteros es conmutativa.
  \item Todo número racional distinto de cero tiene inverso multiplicativo en $\Q$.
  \item El conjunto $\N$ es cerrado para la resta.
\end{opciones}
\begin{solucion}[2.5cm]
a) Falsa: $5 - 3 = 2$ pero $3 - 5 = -2$. b) Verdadera: el inverso de $\frac{p}{q}$ (con
$p \neq 0$) es $\frac{q}{p}$, que también es racional. c) Falsa: $3 - 5 = -2 \notin \N$.
\end{solucion}

\end{preguntas}

\end{document}
